\documentclass[11pt]{article}
\usepackage[margin=1in]{geometry}
\usepackage{amsmath,amssymb}
\usepackage{booktabs}
\usepackage{xcolor}
\usepackage{listings}
\usepackage[colorlinks=true,linkcolor=blue,urlcolor=blue,citecolor=blue]{hyperref}
\usepackage{parskip}

\definecolor{cset}{RGB}{31,119,180}
\definecolor{cfield}{RGB}{44,160,44}
\definecolor{cop}{RGB}{214,39,40}
\definecolor{cgray}{RGB}{120,120,120}

\lstdefinestyle{yaml}{
  basicstyle=\ttfamily\scriptsize,
  frame=single,
  backgroundcolor=\color{gray!7},
  columns=fullflexible,
  breaklines=true,
  keepspaces=true,
  showstringspaces=false,
  commentstyle=\color{cgray}\itshape,
  morecomment=[l]{\#},
}

\newcommand{\code}[1]{\texttt{#1}}
\newcommand{\op}[1]{\textcolor{cop}{\texttt{#1}}}

\title{\textbf{Cardiomyocyte reentry in Plexus}\\[3pt]
\large An excitable active material: electromechanics on tissue material points}
\author{Janelia Research Campus}
\date{\today}

\begin{document}
\maketitle

\begin{abstract}
This note specifies a Plexus model of a cultured cardiomyocyte sheet as an
\textbf{excitable active material}, discretized by \textbf{material points} rather
than segmented cells. We have no cell segmentation, and we do not pretend each node
is a biological cell: the sheet is a fixed grid of \emph{tissue particles}, each
carrying both an electrical state $(u,w)$ and a mechanical state
$(x,v,F,C,m,\mu,\lambda)$. This is closer to a continuum discretization than to a
cell model --- and Plexus expresses both. The design respects the Plexus contract
(Sets / Fields / Operators; the spec is the program, the registry the vocabulary,
the engine the interpreter): every mechanism is a typed operator, and no new
dynamics enter the engine. The electrical kernel reuses the FitzHugh--Nagumo
dynamics already implemented in \textsc{NeuralGraph} (\code{ODEs/Fitzhug\_Nagumo.py});
the mechanics reuse the existing MLS-MPM stack (\code{p2g}, \code{mpm\_grid\_update},
\code{g2p}). Because electrical and mechanical state share one set, the
excitation--contraction and mechano-electric couplings are \emph{local} (per
particle): the up/down \code{aggregate}/\code{broadcast} pair is not needed. The one
methodological constraint that governs everything is that the only external input is
a \emph{sparse, explicit trigger} --- never a learned field that could explain the
data arbitrarily. The rotating reentrant pattern must arise from the excitable
dynamics and the tissue geometry, and must persist after the trigger stops.
\end{abstract}

\section{The tissue as an excitable active material}

We model the cardiomyocyte sheet as a fixed set of \textbf{material points} (tissue
particles), not cells:
\[
  S_{\mathrm{site}}=\{1,\ldots,N\},\qquad N = 100\times 100 = 10{,}000,
\]
arranged on a regular grid. \textbf{The nodes are not cells} --- there is no
segmentation. They are material sites that discretize the active medium; several
sites cover what would biologically be one cardiomyocyte, exactly as in a continuum
finite-element or material-point discretization. The biological claim is therefore
clean: \emph{we model the cardiomyocyte sheet as an excitable active material,
discretized by material particles}.

Each site carries an electrical state and a mechanical state on the \emph{same} row:
\[
  s_i = \big(\underbrace{u_i,\,w_i}_{\text{electrical}},\;
              \underbrace{x_i,\,v_i,\,F_i,\,C_i,\,m_i,\,\mu_i,\,\lambda_i}_{\text{mechanical (MPM)}}\big).
\]
Here $u_i$ is excitation and $w_i$ recovery; $x_i,v_i$ are position and velocity,
$F_i$ the deformation gradient, $C_i$ the affine velocity, $m_i$ mass, and
$\mu_i,\lambda_i$ the Lam\'e parameters.

\paragraph{Material adjacency.}
The sites are joined by a fixed relation
\[
  E_{\mathrm{grid}}\subseteq S_{\mathrm{site}}\times S_{\mathrm{site}},
\]
a four- or eight-neighbour lattice. This is the scaffold for electrical propagation
(the gap-junction analogue). It is a property of the \emph{material}, built once and
never rebuilt from deformed positions.

\paragraph{External trigger.}
We avoid a learnable field that could explain the data arbitrarily. The only
external input is a sparse trigger current applied to a prescribed region
$R\subset S_{\mathrm{site}}$ at discrete times:
\[
  I_i(t)=A\,\mathbf{1}_{i\in R}\,\mathbf{1}_{\,t\bmod T < \Delta T},
\]
with amplitude $A$, period $T$, pulse duration $\Delta T$. The rotating pattern must
arise from the excitable dynamics and the tissue geometry, not from a hidden driving
field.

\paragraph{Excitable dynamics.}
Electrical activity is a FitzHugh--Nagumo operator on the material graph:
\[
  \dot u_i
    = u_i-\frac{u_i^3}{3}-w_i
      + D\!\!\sum_{j:(i,j)\in E_{\mathrm{grid}}}\!\!(u_j-u_i)
      + I_i(t),
  \qquad
  \dot w_i = \varepsilon\,(u_i+a-b\,w_i).
\]
The graph Laplacian propagates excitation between neighbouring sites; $w_i$ gives
refractoriness. This is the minimal mechanism that supports traveling waves, wave
break, and reentry. The single-site core ($D{=}0$) is exactly the dynamics validated
in \textsc{NeuralGraph} (\code{ODEs/Fitzhug\_Nagumo.py}): same
$\dot v=v-v^3/3-w+I_{\mathrm{ext}}$, $\dot w=\varepsilon(v+a-bw)$, Euler stepped,
same pulse protocol (defaults $a{=}0.7$, $b{=}0.8$, $\varepsilon{=}0.18$). Plexus
adds only the spatial coupling over $E_{\mathrm{grid}}$.

\paragraph{Minimal Plexus form (Stage 1).}
\[
  \mathrm{grid\_graph}
  \;\longrightarrow\;
  \mathrm{trigger\_pulse}
  \;\longrightarrow\;
  \mathrm{excitable\_nagumo},
\]
\[
  (u,w)_{t+1}
    = \Phi_{\mathrm{Nagumo}}\!\Big((u,w)_t\,;\,
        E_{\mathrm{grid}},\, I(t),\, D,\,\varepsilon,\,a,\,b\Big).
\]

\section{Electromechanics on one set (local couplings)}

Because the electrical state $(u,w)$ and the mechanical state
$(x,v,F,C,\dots)$ live on the \emph{same} material point, the two
electromechanical couplings are \textbf{local} --- a per-site read/write, with no
parent map to traverse. This is the simplification the material-point choice buys:
the \code{aggregate}/\code{broadcast} pair (needed only to move information between a
cell and its particles) disappears.

\paragraph{Signal to mechanics (local).}
Excitation produces active contraction at the same site:
\[
  f_i^{\mathrm{act}} = \alpha\,g(u_i)\,n_i,
  \qquad
  g(u_i)=\sigma\!\left(\frac{u_i-\theta}{\eta}\right)
  \;\;\text{or}\;\;
  g(u_i)=\mathbf{1}_{\,u_i>\theta},
\]
where $\alpha$ is contractility, $n_i$ a preferred contraction direction, and $g$ a
smooth or thresholded gate. There is no broadcast: $f_i^{\mathrm{act}}$ is written on
site $i$ as an active stress that the MPM solver reads.

\paragraph{MPM mechanics.}
The material sheet is advanced by the MLS-MPM exchange--field--exchange cycle through
the background MPM grid (a mechanical \emph{field}, not a biological one):
\[
  S_{\mathrm{site}}
  \xrightarrow{\;\mathrm{P2G}\;}
  F_{\mathrm{MPM}}
  \xrightarrow{\;\mathrm{grid\_update}\;}
  F_{\mathrm{MPM}}
  \xrightarrow{\;\mathrm{G2P}\;}
  S_{\mathrm{site}},
\]
\[
  (x,v,F,C)_{t+1}
    = \Phi_{\mathrm{MPM}}\!\Big((x,v,F,C)_t\,;\,
        f^{\mathrm{act}},\,\mu,\,\lambda\Big).
\]

\paragraph{Mechanics to signal (local).}
Each site's own deformation gives a local stress/strain summary
\[
  \sigma_i = s(F_i)
\]
(a scalar invariant of the site's deformation gradient --- no fibre average, since
there are no cells to average over). It feeds back into the excitable dynamics as a
mechanosensitive current,
\[
  I_i(t)=I_i^{\mathrm{trigger}}(t)+\beta\,\sigma_i,
\]
or a stress-dependent threshold, $a_i=a_0+\beta\,\sigma_i$. The coupled electrical
equation becomes
\[
  \dot u_i
    = u_i-\frac{u_i^3}{3}-w_i
      + D\!\!\sum_{j:(i,j)\in E_{\mathrm{grid}}}\!\!(u_j-u_i)
      + I_i^{\mathrm{trigger}}(t)
      + \beta\,\sigma_i.
\]

\section{Closed-loop model}

The full electromechanical loop runs entirely on $S_{\mathrm{site}}$:
\[
  u_i
  \longrightarrow f_i^{\mathrm{act}}
  \longrightarrow (x_i,v_i,F_i)
  \longrightarrow \sigma_i
  \longrightarrow u_i,
\]
a composition of typed operators
\[
  \mathrm{trigger}\;\circ\;\mathrm{Nagumo}\;\circ\;
  \mathrm{signal\_to\_force}\;\circ\;\mathrm{MPM}\;\circ\;\mathrm{stress\_to\_signal}.
\]
Stage~1 asks whether an explicit trigger plus excitable graph dynamics produce the
observed rotating curved-string pattern; the later stages ask whether mechanical
feedback stabilizes, bends, or sustains that rotating wave.

\section{Operators: what exists, what is new}

Each row is one operator, classified by the Plexus \emph{kind} it belongs to (what
it changes), with its status in the current registry
(\code{src/plexus/operators/}). All operators act on the single set
\code{tissue\_particle}.

\begin{center}
\renewcommand{\arraystretch}{1.25}
\footnotesize
\begin{tabular}{@{}llll@{}}
\toprule
\textbf{operator} & \textbf{kind} & \textbf{role} & \textbf{status} \\
\midrule
\op{grid\_graph}            & rewire        & fixed material adjacency $E_{\mathrm{grid}}$            & \textbf{new} (fixed-once variant of \code{radius\_graph}) \\
\op{trigger\_pulse}         & lateral (source) & sparse S1/S2 current into $u$ on $R$                 & \textbf{new} (small) \\
\op{excitable\_nagumo}      & lateral       & FHN reaction $+$ graph Laplacian; updates $u,w$         & \textbf{new} (kernel from \textsc{NeuralGraph}) \\
\op{signal\_to\_mpm\_force} & lateral (local) & $u_i\!\to$ active contraction stress on site $i$       & \textbf{new} (local, no broadcast) \\
\op{mls\_mpm\_mechanics}    & exchange ($+$mpm grid) & deforms the sheet; $x,v$ at 2nd order           & \textbf{exists} \\
\op{mpm\_stress\_to\_signal}& lateral (local) & site stress $\sigma_i\!\to$ excitability / threshold   & \textbf{new} (local, no aggregate) \\
\bottomrule
\end{tabular}
\end{center}

The loop uses three kinds --- \code{rewire}, \code{lateral}, \code{exchange} (the MPM
grid is the only field) --- and \emph{not} \code{aggregate}/\code{broadcast}: with
electrical and mechanical state on one material set, every coupling is local.
\code{divide}/\code{die} are unused (no proliferation). Four small new operators are
required; the MPM stack already exists.

\paragraph{On the kind of \code{trigger\_pulse}.}
Conceptually the trigger is a \emph{set-local source}: it injects current into $u$
on a region, not along edges. Plexus has no \code{source} kind yet, so the pragmatic
implementation is \code{class TriggerPulse(Lateral)} with \code{PREDICTION = None},
documented to mutate the auxiliary state $u$ in place (returning \code{\{\}}); a
dedicated \code{source} kind can be promoted later. It honours the selector
\code{tissue\_particle[...]} that scopes it to $R$.

\section{Constraints the contract enforces}

\begin{enumerate}\itemsep4pt
\item \textbf{The nodes are material points, not cells.} No segmentation is claimed.
The set is \code{tissue\_particle} on a grid; the biological object is an excitable
active \emph{material}, and several sites tile one cardiomyocyte. This keeps the
modelling honest about what the data support (a displacement field, not cell
outlines).

\item \textbf{The material adjacency is fixed.} \code{grid\_graph} is built once and
never rebuilt from deformed positions (unlike \code{radius\_graph}, which rewires
each tick). Mechanical deformation does not change $E_{\mathrm{grid}}$. Stretch may
modify conduction \emph{parameters} (via \code{mpm\_stress\_to\_signal}); it must not
rewire the graph.

\item \textbf{$u,w$ are auxiliary state; $x,v$ are integrated.} On one set these
coexist cleanly: \code{excitable\_nagumo} Euler-steps $u,w$ in place (the
\code{heading} precedent in \code{models/base.py}) and returns \code{\{\}}, while
only $x,v$ pass through the engine's second-order integrator (via \code{mpm}). There
is no prediction-order conflict because the engine integrates only $x,v$.

\item \textbf{The trigger must be explicit and sparse.} The only external input is
\code{trigger\_pulse}, scoped by a selector and gated in time. No learned or dense
driving field is permitted, so the reentrant pattern is a consequence of the
dynamics, not an artifact of an over-expressive input.

\item \textbf{The rotor must persist after the trigger stops.} The anti-cheating
metric: if the rotating wave dies once the trigger ends, the pulse was carrying the
pattern. Sustained rotation with the trigger \emph{off} is the criterion for success.
\end{enumerate}

\section{Implementation roadmap}

Each stage ships an \emph{intent} metric --- ``it ran'' is not ``it worked''.

\paragraph{Stage 1a --- planar conduction.}
A fixed $100\times100$ material grid; \code{grid\_graph} built once; a single S1
\code{trigger\_pulse} at one edge; \code{excitable\_nagumo}. The mechanical state is
dormant. \emph{Intent metric:} a clean traveling wavefront with a stable conduction
velocity.

\begin{lstlisting}[style=yaml,caption={Stage 1a spec sketch (electrical only; names match the registry).}]
general: {name: cardio_1a, seed: 0, n_frames: 600, dt: 0.1, boundary: wall}
sets:
  tissue_particle:
    layout: grid
    shape: [100, 100]          # 10000 material sites; state = x, v, u, w, F, C, m, mu, lambda
operators:
  - {op: grid_graph,       at: tissue_particle, neighbours: 8, fixed: true}
  - {op: trigger_pulse,    at: "tissue_particle[x<0.03]", into: u, amp: 0.8, period: 1e9, dur: 1.0}
  - {op: excitable_nagumo, at: tissue_particle, D: 0.1, a: 0.7, b: 0.8, eps: 0.18}
schedule: [trigger_pulse, excitable_nagumo]      # grid_graph built once at init
# stability: dt * D < 0.25 (Laplacian CFL)
\end{lstlisting}

\paragraph{Stage 1b --- rotor via S1--S2.}
\code{trigger\_pulse} fires two pulses: S1 launches a planar wave, S2 fires later on
a partial sub-domain in the refractory tail of S1, so the front curls into a rotor.
\emph{Intent metric:} a sustained rotor \emph{after the trigger ends}, with its
spiral-tip trajectory and rotation period.

\paragraph{Stage 2 --- excitation--contraction coupling.}
Activate the mechanical state on the same sites: \code{signal\_to\_mpm\_force} writes
$\alpha g(u_i)n_i$ locally, \code{mls\_mpm\_mechanics} deforms the sheet.
\emph{Intent metric:} contraction strain locked to the wave phase.

\begin{lstlisting}[style=yaml,caption={Stage 2 spec sketch: add the local force and the MPM cycle on the same set.}]
operators:
  - {op: grid_graph,          at: tissue_particle, neighbours: 8, fixed: true}
  - {op: trigger_pulse,       at: "tissue_particle[...]", into: u, ...}
  - {op: excitable_nagumo,    at: tissue_particle, D: 0.1, a: 0.7, b: 0.8, eps: 0.18}
  - {op: signal_to_mpm_force, at: tissue_particle, alpha: 1.0, theta: 0.5}   # u -> active stress (local)
  - {op: mls_mpm_mechanics,   at: tissue_particle, youngs: 50, n_grid: 128, substeps: 8}
schedule: [trigger_pulse, excitable_nagumo, signal_to_mpm_force, mls_mpm_mechanics]
\end{lstlisting}

\paragraph{Stage 3 --- mechano-electric feedback.}
\code{mpm\_stress\_to\_signal} reads each site's stress and feeds it back as
$+\beta\sigma_i$ (current) or $a_0+\beta\sigma_i$ (threshold). The sign is physical:
stretch-activated channels depolarize, which can be pro- or anti-arrhythmic.
\emph{Intent metric:} whether feedback changes the rotor's drift / anchoring /
breakup relative to Stage~1b.

\paragraph{Stage 4 --- mechanism search.}
Search the coupling parameters $(D,\alpha,\theta,\beta,\dots)$ and, more importantly,
the coupling \emph{family} (current vs.\ threshold; sign of $\beta$) with the bandit
search already used in the prototype. The branches are hypotheses about
mechano-electric feedback, not merely numerical settings.

\section{Provenance and data}

\paragraph{Electrical kernel.} Reused from \textsc{NeuralGraph}
(\code{src/NeuralGraph/ODEs/}): \code{Fitzhug\_Nagumo.py} (single-site FHN generator
$+$ inverse-problem trainer; same equations and rectangular pulse train) and
\code{config\_Fitzhug\_Nagumo.py} (typed config; defaults $a{=}0.7,b{=}0.8,
\varepsilon{=}0.18,dt{=}0.1$). Porting is: lift the per-site FHN update into
\code{excitable\_nagumo}, add the graph-Laplacian coupling over $E_{\mathrm{grid}}$,
and replace the scalar pulse with the selector-scoped \code{trigger\_pulse}.

\paragraph{Real data (the fit target).} The Utrecht cardiomyocyte dataset lives at
\code{.../GraphData/graphs\_data/cardiomyocytes\_real\_data/Cardio\_1/}: a raw
microscope movie (\code{0\_B\_15kPa\_..ome.tif}, $[240,2048,2048]$; an HCM/diseased
movie is also present) and its \code{.derivatives.npy} ($[239,137,137,12]$; channels
$0,1$ are the tracked control-grid coordinates $X(t),Y(t)$ --- the displacement
field). There is no cell segmentation; the tracked grid stands in for material-point
motion. Rendering scripts and three mp4s (raw movie; movie $+$ tracked nodes;
$10\times10$ amplified node trajectories) are in
\code{Plexus/prototype/cardio/} alongside this document.

\end{document}
